% Part: history
% Chapter: biographies
% Section: gerhard-gentzen

\documentclass[../../../include/open-logic-section]{subfiles}

\begin{document}

\olfileid{his}{bio}{gen}

\olsection{Gerhard Gentzen}

\olphoto{gentzen-gerhard}{Gerhard Gentzen}

Gerhard Gentzen 以结构证明论的创立者闻名，尤以创立自然演绎和相继式演算推导系统而著称。他于 1909 年 11 月 24 日出生在德国格赖夫斯瓦尔德。Gerhard 在家接受三年家庭教育后才进入预科学校，当时学业上落后于大多数同学。尽管如此，他依然是一名出色的学生，展现出极高的数学天赋。他兴趣广泛，比如曾为母亲写诗，为学校剧院写剧本。

Gentzen 在格赖夫斯瓦尔德大学开始大学生涯，但后来先后转学至哥廷根、慕尼黑和柏林。他于 1933 年在 Hermann Weyl（赫尔曼·外尔）的指导下获得哥廷根大学博士学位。（其大部分研究由 Paul Bernays（保罗·贝尔奈斯）指导，但 Bernays 遭纳粹解聘。）1934 年，Gentzen 开始担任 希尔伯特的助手。同年，他在论文《逻辑推演研究 I–II》中建立了相继式演算和自然演绎推导系统。他于 1936 年证明了 皮亚诺公理的一致性。

Gentzen 与纳粹的关系颇为复杂。就在其导师 Bernays 被迫离开德国的同时，Gentzen 却加入了纳粹准军事组织冲锋队（SA）的大学支部。和许多德国人一样，他也是纳粹党成员。战争期间，他在空军情报部门担任通信军官。然而，1942 年他因神经衰弱而被免职。目前尚不清楚 Gentzen 是否效忠纳粹党，亦不清楚他加入该党是否是为了确保学术上的成功。

1943 年，Gentzen 获得了布拉格德国大学数学研究所的教职并应聘前往。然而，1945 年布拉格市民爆发了反抗德国占领的起义。苏联军队进驻该市，逮捕了大学里所有的教授。由于身为纳粹组织成员，Gentzen 被押往强制劳改营。他于 1945 年 8 月 4 日在牢房中因营养不良去世，年仅 35 岁。

\begin{reading}
关于 Gentzen 的完整传记，参见 \citet{Menzler-Trott2007}。关于纳粹统治下的数学家的一部有趣著作（其中对 Gentzen 的生平有简要介绍）可参见 \citet{Segal2014}。Gentzen 关于逻辑推演的论文可见德文原版 \citep{Gentzen1935a,Gentzen1935b}。Gentzen 论文的英文译文已由 \citet{Gentzen1969} 汇编为一卷出版，其中也包含一篇生平简介。
\end{reading}

\end{document}